\documentclass[../../main.tex]{subfiles}% 注意这里的文件路径不能用 ./main.tex ,否则用latexmk编译子文件会报错
\graphicspath{{\subfix{./image/}}{\subfix{../../image/}}} % 指定图片引用路径,后续可以直接使用图片文件名
% 如果图片存放在上级目录,则使用 ../../image/ 作为备用路径

% 例如:
% \begin{figure}[H]
% \centering
% \includegraphics[width=0.3\textwidth]{图.png}
% \caption{}
% \label{figure:图}
% \end{figure}
% 注意:上述\label{}一定要放在\caption{}之后,否则引用图片序号会只会显示??.

\begin{document}

\section{调和函数}

\subsection{实例}

\begin{definition}
设 $\Omega\subset \mathbb{R}^n$,如果函数 $u:\Omega\to\mathbb{R}$ 具有二阶连续偏导数且满足 Laplace 方程 $\Delta u = 0$，我们称之为\textbf{调和函数}；如果一个多项式满足 Laplace 方程 $\Delta u = 0$，我们称之为\textbf{调和多项式}。
\end{definition}

\begin{proposition}
设 $\Omega\subset \mathbb{R}^n$,$n\in \mathbb{N}$,则下列函数均为调和函数:
\begin{enumerate}[(1)]
\item 常数函数 $1$ 与线性函数 $x_i$ $(i=1,\cdots,n)$；
\item 二次多项式 $x_i^2-x_j^2$ $(i,j=1,\cdots,n)$；
\item 若 $\Omega\subset\mathbb R^2$ 为开集，视其为复平面区域，$f:\Omega\to\mathbb C$ 解析，则 $u=\Re f$ 与 $v=\Im f$ 调和；
\item 特别地，对任意非负整数 $k$，$z^k=(x+iy)^k$ 的实部 $P_k(x,y)$ 与虚部 $Q_k(x,y)$ 是 $k$ 次调和多项式，且满足递推关系
\begin{align}
P_{k+1}=xP_k-yQ_k,\qquad Q_{k+1}=xQ_k+yP_k. \label{rec:pq}
\end{align}
\item 函数 $e^x\cos y$ 与 $e^x\sin y$ 调和。
\end{enumerate}
\end{proposition}
\begin{proof}
\begin{enumerate}[(1)]
\item 直接计算得
\begin{align*}
\Delta 1=0,\qquad \Delta x_i=0.
\end{align*}

\item 注意到
\begin{align*}
\Delta(x_i^2-x_j^2)=2-2=0.
\end{align*}

\item 由 Cauchy-Riemann 方程
\begin{align*}
\frac{\partial u}{\partial x}=\frac{\partial v}{\partial y},\qquad 
\frac{\partial u}{\partial y}=-\frac{\partial v}{\partial x},
\end{align*}
可得
\begin{align*}
\frac{\partial^2 u}{\partial x^2}+\frac{\partial^2 u}{\partial y^2}
= \frac{\partial}{\partial x}\left(\frac{\partial v}{\partial y}\right)
+\frac{\partial}{\partial y}\left(-\frac{\partial v}{\partial x}\right)
= v_{yx}-v_{xy}=0,
\end{align*}
同理 $\Delta v=0$，故 $u,v$ 调和。

\item 由于 $z^k$ 是解析函数，由 (3) 知其实部 $P_k$ 与虚部 $Q_k$ 调和；递推关系 \eqref{rec:pq} 由 $z^{k+1}=z^k\cdot z$ 直接展开得到。

\item 因为 $e^z=e^x(\cos y+i\sin y)$ 解析，其实部 $e^x\cos y$ 与虚部 $e^x\sin y$ 调和。
\end{enumerate}
\end{proof}

\begin{theorem}
假设 $u(x)$ 是一个 $\mathbb{R}^n$ 上的调和函数，则
\begin{enumerate}[(1)]
\item $u(\lambda x)$ 是一个调和函数，其中 $\lambda$ 是任一实数；
\item $u(x+x_0)$ 是一个调和函数，其中 $x_0\in\mathbb{R}^n$ 固定；
\item $u(Ox)$ 是一个调和函数，其中 $O:\mathbb{R}^n\to\mathbb{R}^n$ 是一个正交变换。
\end{enumerate}
\end{theorem}
\begin{remark}
此定理说明在伸缩变换、平移变换和正交变换下调和函数仍变为调和函数。
\end{remark}



\subsection{平均值公式}

\begin{theorem}[调和函数平均值公式]\label{theorem:调和函数平均值公式}
设 $\Omega\subset \mathbb{R}^n$, $u\in C^2(\Omega)$ 是 $\Omega$ 上的调和函数，则对于任意的球 $B(x,r)\subset\Omega$，有
\begin{align}
u(x) = \fint_{\partial B(x,r)} u(y)\,\mathrm{d}S(y) = \fint_{B(x,r)} u(y)\,\mathrm{d}y, \label{eq:mean-value}
\end{align}
其中积分号 $\fint$ 表示求平均值，即
\begin{gather*}
\fint_{\partial B(x,r)} u(y)\,\mathrm{d}S(y) = \frac{1}{n\alpha(n)r^{n-1}} \int_{\partial B(x,r)} u(y)\,\mathrm{d}S(y), \\
\fint_{B(x,r)} u(y)\,\mathrm{d}y = \frac{1}{\alpha(n)r^n} \int_{B(x,r)} u(y)\,\mathrm{d}y,
\end{gather*}
这里 $\alpha(n)=\frac{2\pi^{\frac{n}{2}}}{n\Gamma\left(\frac{n}{2}\right)}$ 表示 $\mathbb{R}^n$ 上单位球的体积。
\end{theorem}
\begin{proof}
令
\begin{align*}
\phi(r) = \fint_{\partial B(x,r)} u(y)\,\mathrm{d}S(y).
\end{align*}
作平移和伸缩变换，则
\begin{align*}
\phi(r) = \fint_{\partial B(0,1)} u(x+rz)\,\mathrm{d}S(z).
\end{align*}
对 $r$ 求微商，由\hyperref[Basis of Analytics-theorem:含参量积分的可微性1]{含参量积分的可微性}和\hyperref[Basis of Analytics-theorem:Green第一恒等式]{Green第一恒等式}我们得到
\begin{align*}
\phi'(r) &= \fint_{\partial B(0,1)} D u(x+rz)\cdot z\,\mathrm{d}S(z)
= \fint_{\partial B(x,r)} D u(y)\cdot \frac{y-x}{r}\frac{1}{r^{n-1}}\,\mathrm{d}S(y) \\
&= \frac{1}{r^{n-1}}\fint_{\partial B(x,r)} D u(y)\cdot n\,\mathrm{d}S(y) = \frac{1}{r^{n-1}}\fint_{\partial B(x,r)} \frac{\partial u(y)}{\partial n}\,\mathrm{d}S(y) \\
&= \frac{1}{r^{n-1}}\fint_{B(x,r)} \Delta u(y)\,\mathrm{d}y = 0,
\end{align*}
于是 $\phi(r)$ 是一个常数。由 $u(x)$ 的连续性，从而
\begin{align*}
u(x) = \lim_{t\to 0^+} \phi(t) = \phi(r),
\end{align*}
即
\begin{align}
u(x) = \fint_{\partial B(x,r)} u(y)\,\mathrm{d}S(y). \label{eq:mean-sphere}
\end{align}

注意到
\begin{align*}
\int_{B(x,r)} u(y)\,\mathrm{d}y = \int_0^r \int_{\partial B(x,t)} u(y)\,\mathrm{d}S(y)\,\mathrm{d}t,
\end{align*}
再利用等式 \eqref{eq:mean-sphere}，则
\begin{align*}
\int_{B(x,r)} u(y)\,\mathrm{d}y = u(x)\int_0^r n\alpha(n)t^{n-1}\,\mathrm{d}t = \alpha(n)r^n u(x).
\end{align*}
于是我们完成定理的证明.
\end{proof}

\begin{theorem}\label{theorem:调和函数平均值定理逆命题}
设 $\Omega\subset \mathbb{R}^n$, $u\in C^2(\Omega)$ 满足，对于任意 $B(x,r)\subset\Omega$，
\begin{align*}
u(x) = \fint_{\partial B(x,r)} u(y)\,\mathrm{d}S(y),
\end{align*}
则 $u$ 是调和函数。
\end{theorem}
\begin{proof}
任取$x\in\Omega$，固定$x$,对任意球 $B(x,r)\subset\Omega$，由\hyperref[theorem:调和函数平均值公式]{调和函数平均值公式}知
\begin{align*}
\phi(r) = \fint_{\partial B(x,r)} u(y)\,\mathrm{d}S(y)= \int_{\partial B(0,1)} u(x+rz)\,\mathrm{d}S(z)
\end{align*}
是一个常数.因此
$\phi'(r) = 0.$
再由\hyperref[Basis of Analytics-theorem:含参量积分的可微性1]{含参量积分的可微性}和\hyperref[Basis of Analytics-theorem:Green第一恒等式]{Green第一恒等式}我们得到
\begin{align*}
\phi'(r) &= \fint_{\partial B(0,1)} D u(x+rz)\cdot z\,\mathrm{d}S(z)
= \fint_{\partial B(x,r)} D u(y)\cdot \frac{y-x}{r}\frac{1}{r^{n-1}}\,\mathrm{d}S(y) \\
&= \frac{1}{r^{n-1}}\fint_{\partial B(x,r)} D u(y)\cdot n\,\mathrm{d}S(y) = \frac{1}{r^{n-1}}\fint_{\partial B(x,r)} \frac{\partial u(y)}{\partial n}\,\mathrm{d}S(y) \\
&= \frac{1}{r^{n-1}}\fint_{B(x,r)} \Delta u(y)\,\mathrm{d}y = 0,
\end{align*}
于是对任意球 $B(x,r)\subset\Omega$,都有
\begin{align}\label{eq::HOUFHwie22312}
\int_{B(x,r)} \Delta u(y)\,\mathrm{d}y = 0.
\end{align}
若$\Delta u\left( x \right) \ne 0$, 不妨设$\Delta u\left( x \right) >0$. 因为$\Delta u$在$x$处连续, 所以$\exists \delta >0$, 使得
\begin{align*}
\Delta u\left( y \right) >\frac{\Delta u\left( x \right)}{2}>0,\quad \forall y\in B\left( x,\delta \right).
\end{align*}
从而
\begin{align*}
\int_{B\left( x,\delta \right)}\Delta u\left( y \right)\,\mathrm{d}y>\int_{B\left( x,\delta \right)}\frac{\Delta u\left( x \right)}{2}\,\mathrm{d}y=\frac{\Delta u\left( x \right)}{2}\cdot \alpha \left( n \right)\delta ^n>0,
\end{align*}
这与\eqref{eq::HOUFHwie22312}式矛盾!故$\Delta u(x)=0.$(也可由积分中值定理再令$r\to 0^+$得到)\end{proof}

\begin{theorem}\label{theorem:theorem:调和函数平均值定理逆命题-u弱化成连续}
设 $\Omega\subset \mathbb{R}^n$, $u\in C(\Omega)$ 满足平均值公式，即对于任意 $B(x,r)\subset\Omega$，有
\begin{align*}
u(x) = \fint_{\partial B(x,r)} u(y)\,\mathrm{d}S(y),
\end{align*}
则 $u$ 是调和函数且 $u\in C^\infty(\Omega)$。
\end{theorem}
\begin{remark}
定理中 $u\in C(\Omega)$ 的假设还可以降低，我们在这里就不展开讨论了。
\end{remark}
\begin{proof}
设 $\eta:\mathbb{R}^n\to\mathbb{R}$ 是一个光滑化子(\hyperref[Real Analysis-definition:标准磨光核]{标准磨光核})，即径向对称的非负函数 $\eta\in C_0^\infty(\mathbb{R}^n)$ 满足：
\begin{enumerate}[(1)]
\item $\eta(x)=\zeta(|x|)$，这里 $\zeta:[0,+\infty)\to[0,+\infty)$ 是一个非负函数；
\item 其支集 $\operatorname{spt}\eta\subset B(0,1)$；
\item $\int_{\mathbb{R}^n}\eta(x)\,\mathrm{d}x = \int_{B(0,1)}\eta(y)\,\mathrm{d}y = 1$.
\end{enumerate}
对于 $\forall \varepsilon>0$，定义
\begin{align*}
\eta_\varepsilon(y) = \frac{1}{\varepsilon^n}\eta\left(\frac{y}{\varepsilon}\right),
\end{align*}
于是 $\operatorname{spt}\eta_\varepsilon\subset B(0,\varepsilon)$，且
\begin{align*}
\int_{\mathbb{R}^n}\eta_\varepsilon(y)\,\mathrm{d}y
= \frac{1}{\varepsilon^n}\int_{B(0,\varepsilon)} \eta\left(\frac{y}{\varepsilon}\right)\,\mathrm{d}y 
= \int_{B(0,1)}\eta(y)\,\mathrm{d}y = 1.
\end{align*}
利用球坐标变换,再结合上式我们得到
\begin{align}
\int_{\mathbb{R} ^n}{\eta _{\varepsilon}(y)\,\mathrm{d}y}&=\frac{1}{\varepsilon ^n}\int_{B(0,\varepsilon )}{\zeta \left( \frac{\left| y \right|}{\varepsilon} \right) \,\mathrm{d}y}=\frac{1}{\varepsilon ^n}\int_0^{\varepsilon}{\zeta \left( \frac{r}{\varepsilon} \right) \,\left( \int_{S^{n-1}}{1\,}\mathrm{d}\sigma \left( \theta \right) \right) \mathrm{d}r} \nonumber 
\\
&=\frac{1}{\varepsilon ^n}\int_0^{\varepsilon}{\zeta \left( \frac{r}{\varepsilon} \right)}\cdot n\alpha (n)r^{n-1}\,\mathrm{d}r=1.\label{eq::hjijeioefjwjfiowejiefi2324}
\end{align}
对于充分小的 $\varepsilon>0$，在区域 $\Omega_\varepsilon = \{x\in\Omega \mid \operatorname{dist}(x,\partial\Omega)>\varepsilon\}$ 上定义
\begin{align*}
u_\varepsilon(x) = \int_\Omega \eta_\varepsilon(x-y)u(y)\,\mathrm{d}y.
\end{align*}
由于 $\eta$ 无穷次可微且具有紧支集，因此由\refthe{Real Analysis-theorem:实分析--卷积磨光性质}知$u_\varepsilon\in C^\infty(\Omega_\varepsilon)$。对于任意$x\in\Omega_\varepsilon$，我们利用题设已知的平均值公式和\eqref{eq::hjijeioefjwjfiowejiefi2324}式得
\begin{align*}
u_\varepsilon(x) &= \frac{1}{\varepsilon^n}\int_{B(x,\varepsilon)} \eta\left(\frac{x-y}{\varepsilon}\right)u(y)\,\mathrm{d}y = \frac{1}{\varepsilon^n}\int_0^\varepsilon \zeta\left(\frac{r}{\varepsilon}\right) \int_{\partial B(x,r)} u(y)\,\mathrm{d}S(y)\,\mathrm{d}r \\
&= \frac{1}{\varepsilon^n}\int_0^\varepsilon \zeta\left(\frac{r}{\varepsilon}\right) u(x)\cdot n\alpha(n)r^{n-1}\,\mathrm{d}r 
= u(x).
\end{align*}
因而对充分小的 $\varepsilon>0$都有$u\in C^\infty(\Omega_\varepsilon)$。对$\forall x_0\in \Omega$,取$d = \frac{\operatorname{dist}(x_0,\partial\Omega)}{2}>0$,则由$u\in C^{\infty}(\Omega_d)$知$u$在$x = x_0$处无穷次可微.故由$x_0$的任意性知$u\in C^\infty(\Omega)$.
再利用\refthe{theorem:调和函数平均值定理逆命题}即可得证.
\end{proof}

\begin{theorem}[Harnack不等式]\label{theorem:Harnack不等式}
对于 $\Omega\subset \mathbb{R}^n$ 上的任何连通紧子集 $V$，存在一个仅与距离函数
\begin{align*}
\operatorname{dist}(V,\partial\Omega) = \min_{x\in V,\, y\in\partial\Omega} |x-y|
\end{align*}
和维数 $n$ 有关的正常数 $C$，使得
\begin{align*}
\sup_V u \leqslant C \inf_V u,
\end{align*}
其中 $u$ 是 $\Omega$ 上的任意非负调和函数。特别地，对任意 $x,y\in V$，
\begin{align*}
\frac{1}{C}u(y) \leqslant u(x) \leqslant C u(y).
\end{align*}
\end{theorem}
\begin{proof}
取 $r = \frac{1}{4}\operatorname{dist}(V,\partial\Omega)$。
首先考虑 $x,y\in V$，$|x-y|\leqslant r$。于是 $B(y,r)\subset B(x,2r)$，从而由\hyperref[theorem:调和函数平均值公式]{调和函数平均值公式}和$u$的非负性可得
\begin{align*}
u(x)&=\fint_{B\left( x,2r \right)}{u\left( z \right) \mathrm{d}z}=\frac{1}{\left| B\left( x,2r \right) \right|}\int_{B\left( x,2r \right)}{u\left( z \right) \mathrm{d}z}
\\
&=\frac{1}{2^n\left| B\left( y,r \right) \right|}\int_{B\left( x,2r \right)}{u\left( z \right) \mathrm{d}z}\geqslant \frac{1}{2^n\left| B\left( y,r \right) \right|}\int_{B\left( y,r \right)}{u\left( z \right) \mathrm{d}z}
\\
&= \frac{1}{2^n}u\left( y \right) .
\end{align*}
由于 $V$ 是连通的紧集，我们可用一串有限多个半径为 $r$ 的球 $\{B_i\}_{i=1}^N$ 覆盖$V$，而且满足 $B_i\cap B_{i-1}\neq\varnothing$ $(i=2,\cdots,N)$。任取$x_i\in B_i\cap B_{i-1}(i=2,\cdots,N)$,则$|x_i-x_{i-1}|\leqslant r(i=2,\cdots,N)$.
对$\forall x,y\in V$,则由$\{B_i\}_{i=1}^N$ 覆盖$V$知存在$x_k,x_l$满足$|x-x_k|\leqslant r,|y-x_l|\leqslant r$.不妨设$k\leqslant l$,再由上式可得
\begin{align}
u\left( x \right) \geqslant \frac{1}{2^n}u\left( x_k \right) \geqslant \frac{1}{2^n}\left( \frac{1}{2^n}u\left( x_{k+1} \right) \right) =\frac{1}{2^{2n}}u\left( x_{k+1} \right) \geqslant \cdots \geqslant \frac{1}{2^{\left( l-k+1 \right) n}}u\left( x_l \right) \geqslant \frac{1}{2^{\left( l-k+2 \right) n}}u\left( y \right) \geqslant \frac{1}{2^{Nn}}u\left( y \right) .\label{eq::hiowejfiowjio2389049023u8023}
\end{align}
进而
\begin{align}
u(y)\leqslant 2^{nN}u(x),\quad \forall x,y\in V.\label{eq::hiowejfiowjio2389049023u8023-1}
\end{align}
再分别对$y,x$取上、下确界即得
\begin{align*}
\sup_V u \leqslant C \inf_V u.
\end{align*}
特别地,结合\eqref{eq::hiowejfiowjio2389049023u8023}\eqref{eq::hiowejfiowjio2389049023u8023-1}两式也有
\begin{align*}
\frac{1}{2^{nN}}u(y)\leqslant u(x)\leqslant 2^{nN}u(y),\quad \forall x,y\in V.
\end{align*}
\end{proof}

\begin{lemma}
假设 $u$ 是 $B(0,R)\subset \mathbb{R}^n$ 上的非负调和函数。
利用Poisson公式证明：
\begin{align*}
R^{n-2}\frac{R-|x|}{(R+|x|)^{n-1}}u(0) \leqslant u(x) \leqslant R^{n-2}\frac{R+|x|}{(R-|x|)^{n-1}}u(0);
\end{align*}
\end{lemma}
\begin{proof}
\end{proof}

\begin{theorem}[Harnack不等式]\label{theorem:Harnack不等式1}
假设 $u$ 是半径为 $R$ 的球 $B_R\subset \mathbb{R}^n$ 上的任意非负调和函数，则对任意 $r\in(0,R)$，成立不等式
\begin{align*}
\sup_{B_r} u \leqslant \left(\frac{R+r}{R-r}\right)^n \inf_{B_r} u,
\end{align*}
其中 $B_r$ 是以 $r$ 为半径，与 $B_R$ 同心的球。
\end{theorem}
\begin{proof}
\end{proof}

\begin{theorem}[强极值原理]\label{theorem:PED-强极值原理}
假设 $\Omega$ 是 $\mathbb{R}^n$ 上的有界开集，$u\in C^2(\Omega)\cap C(\overline{\Omega})$ 是 $\Omega$ 上的调和函数，则
\begin{enumerate}[(1)]
\item $u(x)$ 在 $\overline{\Omega}$ 上的最大（小）值一定在边界 $\partial\Omega$ 上达到，即
\begin{align*}
\max_{\overline{\Omega }} u=\max_{\partial \Omega} u,\quad \min_{\overline{\Omega }} u=\min_{\partial \Omega} u;
\end{align*}
\item 如果 $\Omega$ 是连通的，且存在 $x_0\in\Omega$ 使得调和函数 $u(x)$ 在 $x_0$ 点达到 $u(x)$ 在 $\overline{\Omega}$ 上的最大（小）值，则 $u$ 在 $\overline{\Omega}$ 上是常数。
\end{enumerate}
\end{theorem}
\begin{remark}
如果只有定理中结论 (1) 成立，我们通常称之为\textbf{弱极值原理}。
\end{remark}
\begin{proof}
仅就最大值的情形证明。
\begin{enumerate}[(1)]
\item 记
\begin{align*}
M = \max_{\overline{\Omega}} u(x).
\end{align*}
如果 $u(x)$ 在边界 $\partial\Omega$ 上的某点达到 $M$，结论 (1) 自然成立。如果 $u(x)$ 在 $\Omega$ 上的内点 $x_0\in\Omega$ 达到 $M$，即 $u(x_0)=M$，我们证明在 $\Omega$ 的一个包含 $x_0$ 的连通分支 $\Omega_1$ 上，调和函数 $u(x)$ 恒为常数 $M$。

固定 $x_1\in\Omega_1$，则由\rrefpro{Topology-proposition:Rn中开集的连通分支必是开集}{Topology-proposition:Rn中开集的连通分支必是开集-1}存在连续映射$\gamma:[0,1]\to\Omega_1$ 连接 $x_0$ 和 $x_1$ 两点，也就是说，路径 $\gamma(t)$ 关于 $t$ 连续且 $\gamma(0)=x_0$ 和 $\gamma(1)=x_1$。定义
\begin{align*}
l = \sup\{t\in[0,1] \mid u[\gamma(t)] = M\},
\end{align*}
我们将证明 $l=1$，从而得到 $u(x_1)=u[\gamma(1)]=M$。

我们用反证法证明 $l=1$。假设 $l<1$，记 $x_l=\gamma(l)$，由函数 $u[\gamma(t)]$ 的连续性，我们有 $u(x_l)=M$。由于 $x_l$ 是内点，因此存在 $B(x_l,r_1)\subset\Omega_1$，在 $B(x_l,r_1)$ 应用\hyperref[theorem:调和函数平均值公式]{调和函数平均值公式}并注意到 $u(x_l)=M$可得
\begin{align}
u(x_l)=\fint_{B\left( x_l,r_1 \right)}{u\left( x \right) \mathrm{d}x}=M\Longrightarrow \int_{B\left( x_l,r_1 \right)}{u\left( x \right) \mathrm{d}x}=M|B(x_l,r_1)|\label{eq::fjisjfiwJGIOJE234}
\end{align}
对$\forall x\in B(x_l,r_1)$, 若$u(x) < M$, 则由$u$的连续性知$\exists \delta > 0$, 使得
\begin{align*}
u(x) < M, \quad \forall x\in B(x,\delta) \Longrightarrow \int_{B(x,\delta)} u(x) \mathrm{d}x < M\left|B(x,\delta)\right|.
\end{align*}
从而
\begin{align*}
\int_{B(x_l,r_1)} u(x) \mathrm{d}x &= \int_{B(x,\delta)} u(x) \mathrm{d}x + \int_{B(x_l,r_1) \setminus B(x,\delta)} u(x) \mathrm{d}x \\
&< M\left|B(x,\delta)\right| + M\left|B(x_l,r_1) \setminus B(x,\delta)\right| 
= M\left|B(x_l,r_1)\right|,
\end{align*}
这与\eqref{eq::fjisjfiwJGIOJE234}式矛盾!
故在$B(x_l,r_1)$ 上函数 $u(x)=M$。注意到 $\gamma(t)$ 在 $l$ 处的连续性，存在充分小的 $\varepsilon>0$，使得 $\gamma$ 在区间 $[l-\varepsilon,l+\varepsilon]$ 上的象集包含于 $B(x_l,r_1)$。于是 $u[\gamma(l+\varepsilon)]=M$。这与 $l$ 是上确界的定义矛盾。因此假设 $l<1$ 不成立，从而 $l=1$。

这样我们证明了在 $\Omega$ 的包含 $x_0$ 的连通分支 $\Omega_1$ 上，$u(x)$ 恒为常数 $M$。由于函数 $u(x)$ 的连续性，$u(x)$ 在 $\Omega_1$ 的边界 $\partial\Omega_1$ 上也恒为常数 $M$。而由\refpro{Topology-proposition:Rn中连通分支的边界包含于原集合的边界}知$\partial\Omega_1$ 是 $\Omega$ 的边界 $\partial\Omega$ 的一部分，从而得到定理的结论 (1) 成立。

\item 当 $\Omega$ 连通时，$\Omega$ 只有一个连通分支，从而 $\Omega_1=\Omega$，由(1)的证明可知定理的结论 (2) 成立。至此我们完成定理的证明。
\end{enumerate}
\end{proof}

\begin{corollary}
假设 $\Omega$ 是 $\mathbb{R}^n$ 上的有界开集且 $g\in C(\partial\Omega)$ 和 $f\in C(\Omega)$，则 Dirichlet 问题 \eqref{eq:dirichlet} 
\begin{align}
\begin{cases}
-\Delta u = f, & x\in\Omega, \\
u = g, & x\in\partial\Omega
\end{cases} \label{eq:dirichlet}
\end{align}
最多存在一个解 $u\in C^2(\Omega)\cap C(\overline{\Omega})$。
\end{corollary}
\begin{proof}
假如存在两个解 $u_1$ 和 $u_2$，并令 $w=u_1-u_2$。由于 Dirichlet 问题 \eqref{eq:dirichlet} 是线性的，故 $w$ 满足齐次边值问题
\begin{align*}
\begin{cases}
-\Delta w = 0, & x\in\Omega, \\
w = 0, & x\in\partial\Omega.
\end{cases}
\end{align*}
我们需要证明 $w\equiv 0$。利用\hyperref[theorem:PED-强极值原理]{极值原理}我们得到，对于 $x\in\Omega$，$w(x)\leqslant 0$。注意到 $-w$ 满足同样的齐次边值问题，于是对于 $x\in\Omega$，$w(x)\geqslant 0$。因而 $w\equiv 0$。至此我们完成推论证明.
\end{proof}

\begin{theorem}\label{theorem:调和函数微分的积分估计}
假设 $u$ 是$\mathbb{R}^n$ 上的有界开集$\Omega$ 上的调和函数，则对于任意球 $B(x,r)\subset\Omega$，任意阶数为 $k$ 的多重指标 $\alpha$，估计
\begin{align}
|D^\alpha u(x)| \leqslant \frac{C_k}{r^{n+k}} \int_{B(x,r)} |u(y)|\,\mathrm{d}y \label{eq:der-est}
\end{align}
成立，其中
\begin{align*}
C_0 = \frac{1}{\alpha(n)},\qquad C_k = \frac{(n+k)^{n+k}(n+1)^k}{\alpha(n)(n+1)^{n+1}} \quad (k=1,2,\cdots). 
\end{align*}
\end{theorem}
\begin{proof}
我们对 $k$ 利用数学归纳法证明公式 \eqref{eq:der-est}。

当 $k=0$ 时，公式 \eqref{eq:der-est}是平均值公式 \eqref{eq:mean-value} 的直接推论；

当 $k=1$ 时，我们注意到 $u_{x_i}$ $(i=1,2,\cdots,n)$ 也是调和函数，于是在球 $B(x,s)$ $(0<s\leqslant r)$ 上利用\hyperref[theorem:调和函数平均值公式]{调和函数平均值公式}和\hyperref[theorem:Green恒等式-3]{Green恒等式}得
\begin{align*}
|u_{x_i}(x)| &= \left| \fint_{B(x,s)} u_{x_i}(y)\,\mathrm{d}y \right| = \left| \frac{1}{\alpha(n)s^n} \int_{\partial B(x,s)} u(y)n_i(y)\,\mathrm{d}S(y) \right|,
\end{align*}
从而有
\begin{align*}
\alpha(n)s^n |u_{x_i}(x)| \leqslant \int_{\partial B(x,s)} |u(y)|\,\mathrm{d}S(y).
\end{align*}
对上式两端在 $(0,r)$ 上积分我们得到
\begin{align*}
|u_{x_i}(x)| \int_0^r \alpha(n)s^n\,\mathrm{d}s \leqslant \int_0^r \mathrm{d}s \int_{\partial B(x,s)} |u(y)|\,\mathrm{d}S(y) = \int_{B(x,r)} |u(y)|\,\mathrm{d}y,
\end{align*}
因而
\begin{align}
|u_{x_i}(x)| \leqslant \frac{n+1}{\alpha(n)r^{n+1}} \int_{B(x,r)} |u(y)|\,\mathrm{d}y. \label{eq:der-k1}
\end{align}
至此我们完成当 $k=1$ 时公式 \eqref{eq:der-est} 的证明。

现在假设 $k\geqslant 2$ 且对于 $\Omega$ 上的所有球和 $k-1$ 阶的多重指标，公式 \eqref{eq:der-est}成立。固定 $B(x,r)\subset\Omega$。令 $\alpha$ 为一个 $k$ 阶多重指标，由\eqref{eq:sfiuhsu-11}式知必然存在 $i\in\{1,2,\cdots,n\}$ 和一个 $k-1$ 阶的多重指标 $\beta$，使得 $D^\alpha u=(D^\beta u)_{x_i}$。令 $v=D^\beta u$。由于 $v$ 是调和函数，对于任意 $0<t<1$，在球 $B(x,tr)$ 上应用不等式 \eqref{eq:der-k1} 有
\begin{align*}
|D^\alpha u(x)| \leqslant \frac{n+1}{\alpha(n)(tr)^{n+1}} \int_{B(x,tr)} |D^\beta u(y)|\,\mathrm{d}y.
\end{align*}
注意到 $y\in B(x,tr)$ 蕴涵着 $B(y,(1-t)r)\subset B(x,r)\subset\Omega$。由归纳假设知,对任意 $0<t<1$，有
\begin{align*}
|D^\beta u(y)| \leqslant \frac{C_{k-1}}{[(1-t)r]^{n+k-1}} \int_{B(y,(1-t)r)} |u(z)|\,\mathrm{d}z \leqslant \frac{C_{k-1}}{[(1-t)r]^{n+k-1}} \int_{B(x,r)} |u(z)|\,\mathrm{d}z.
\end{align*}
综上，我们得到
\begin{align*}
|D^\alpha u(x)| \leqslant \frac{n+1}{\alpha(n)(tr)^{n+1}} \alpha(n)(tr)^n \max_{y\in B(x,tr)} |D^\beta u(y)| \leqslant \frac{(n+1)C_{k-1}}{r^{n+k} t(1-t)^{n+k-1}} \int_{B(x,r)} |u(z)|\,\mathrm{d}z.
\end{align*}
又易证函数 $f(t)=t(1-t)^{n+k-1}$ 在 $t=\frac{1}{n+k}$ 时达到在区间 $(0,1)$ 上的最大值，故在上式中取 $t=\frac{1}{n+k}$，我们得到
\begin{align*}
|D^\alpha u(x)| \leqslant \frac{C_k}{r^{n+k}} \int_{B(x,r)} |u(y)|\,\mathrm{d}y,
\end{align*}
其中
\begin{align*}
C_k = (n+1)C_{k-1} \frac{(n+k)^{n+k}}{(n+k-1)^{n+k-1}}.
\end{align*}
记
\begin{align*}
B_k = \frac{C_k}{(n+k)^{n+k}}, \quad k=1,2,\cdots,
\end{align*}
则我们有递推关系式
\begin{align*}
B_k = (n+1)B_{k-1}, \quad k=2,3,\cdots
\end{align*}
从而
\begin{align*}
B_k = (n+1)^{k-1}B_1, \quad k=2,3,\cdots
\end{align*}
注意到 $C_1=\frac{n+1}{\alpha(n)}$ 和 $B_1=\frac{C_1}{(n+1)^{n+1}}$，我们得到
\begin{align*}
C_k = B_k(n+k)^{n+k} = \frac{(n+k)^{n+k}(n+1)^k}{\alpha(n)(n+1)^{n+1}} \quad (k=1,2,\cdots).
\end{align*}
这就完成当 $|\alpha|=k$ 时公式 \eqref{eq:der-est}的证明，从而由数学归纳法完成定理的证明.
\end{proof}

\begin{theorem}[Liouville定理]\label{theorem:PDE-Liouville定理}
假设 $u$ 是 $\mathbb{R}^n$ 上的有界调和函数，则 $u$ 是常数。
\end{theorem}
\begin{proof}
设 $|u|\leqslant M$。固定 $x\in\mathbb{R}^n$。对任意 $r>0$，在球 $B(x,r)$ 上利用\refthe{theorem:调和函数微分的积分估计}当 $k=1$ 时的结论，则
\begin{align*}
|Du(x)| \leqslant \frac{nC_1}{r^{n+1}} \int_{B(x,r)} |u(y)|\,\mathrm{d}y \leqslant \frac{nC_1\alpha(n)}{r}M.
\end{align*}
令 $r\to+\infty$，我们有
\begin{align*}
Du(x)\equiv 0,\quad x\in\mathbb{R}^n.
\end{align*}
因此由\refpro{Basis of Analytics-proposition:连通开集上梯度为0就是常数}知$u$ 是常数。至此我们完成定理的证明.
\end{proof}

\begin{definition}[上有界和下有界]
\begin{enumerate}[(1)]
\item $u$是$\mathbb{R}^n$的\textbf{上有界}函数是指存在一个常数 $M$，使得对于任意 $x\in\mathbb{R}^n$，有 $u(x)\leqslant M$ 成立.

\item $u$是$\mathbb{R}^n$的\textbf{下有界}函数是指存在一个常数$M$，使得对于任意 $x\in\mathbb{R}^n$，有 $u(x)\geqslant M$ 成立.
\end{enumerate}
\end{definition}

\begin{theorem}
假设 $u$ 是 $\mathbb{R}^n$ 的上有界（或下有界）调和函数，则 $u$ 是一个常数。
\end{theorem}
\begin{proof}
\end{proof}

\begin{theorem}
假设 $u$ 是 $\Omega$ 上的调和函数，则 $u$ 是 $\Omega$ 上的解析函数。
\end{theorem}
\begin{proof}
固定 $x_0\in\Omega$。我们要证 $u(x)$ 在 $x_0$ 的某个邻域内可表示为一个收敛的幂级数。取 $r=\frac{1}{4}\operatorname{dist}(x_0,\partial\Omega)$，并记
\begin{align*}
A = \frac{1}{\alpha(n)(n+1)^{n+1}r^n} \int_{B(x_0,2r)} |u(y)|\,\mathrm{d}y.
\end{align*}
对任意 $x\in B(x_0,r)$，则 $B(x,r)\subset B(x_0,2r)\subset\Omega$。由\refthe{theorem:调和函数微分的积分估计}得到
\begin{align*}
|D^\alpha u(x)| \leqslant \frac{A}{r^N}(n+N)^{n+N}(n+1)^N,
\end{align*}
其中 $|\alpha|=N$。利用 \hyperref[Basis of Analytics-theorem:Stirling公式]{Stirling公式}可得
\begin{align*}
\lim_{N\to\infty} \frac{N^{N+\frac{1}{2}}}{N!e^N} = \frac{1}{(2\pi)^{\frac{1}{2}}} < 1,
\end{align*}
从而当 $N$ 充分大时，有
\begin{align*}
(n+N)^{n+N} = \left(1+\frac{n}{N}\right)^N (N+n)^n N^N 
\leqslant e^n (2N)^n N^N \leqslant 2^n e^n N!e^N N^{n-\frac{1}{2}}.
\end{align*}
我们将要证明，当
$|x-x_0| \leqslant \frac{r}{2n(n+1)e}$
时，$u(x)$ 在 $x_0$ 处的 Taylor 级数
\begin{align*}
\sum_{|\alpha|=0}^{\infty} \frac{D^\alpha u(x_0)}{\alpha!}(x-x_0)^\alpha
\end{align*}
收敛到 $u(x)$。

实际上，对充分大的 $N$，存在 $t=t(x)\in(0,1)$ 使得 $u(x)$ 在 $x_0$ 处的 Taylor 级数的余项可以表示为
\begin{align*}
R_N(x) = u(x) - \sum_{k=0}^{N-1}\sum_{|\alpha|=k} \frac{D^\alpha u(x_0)}{\alpha!}(x-x_0)^\alpha 
= \sum_{|\alpha|=N} \frac{D^\alpha u[x_0+t(x-x_0)]}{\alpha!}(x-x_0)^\alpha.
\end{align*}
利用\hyperref[Basis of Analytics-theorem:组合公式-二项式定理]{二项式定理}知
\begin{align*}
\sum_{|\alpha|=N}\frac{N!}{\alpha!} = n^N,
\end{align*}
我们得到，当 $N$ 充分大时，
\begin{align*}
|R_N(x)| &\leqslant A\sum_{|\alpha|=N} \frac{1}{\alpha!}(n+N)^{n+N}\left[\frac{(n+1)|x-x_0|}{r}\right]^N \\
&\leqslant 2^n e^n A \left(\sum_{|\alpha|=N}\frac{N!}{\alpha!}\right) e^N N^{n-\frac{1}{2}} \left[\frac{(n+1)|x-x_0|}{r}\right]^N \\
&\leqslant 2^n e^n A \left[\frac{n(n+1)e|x-x_0|}{r}\right]^N N^{n-\frac{1}{2}} \leqslant 2^n e^n A \frac{N^{n-\frac{1}{2}}}{2^N},
\end{align*}
从而 $u(x)$ 在 $x_0$ 处的 Taylor 级数的余项 $R_N(x)$ 在 $x_0$ 的上述邻域内一致收敛到 $0$。因此 $u(x)$ 在 $x_0$ 的某个邻域内可表示为一个收敛的幂级数。至此我们完成定理的证明.
\end{proof}

\begin{theorem}
假设 $u$ 是 $\mathbb{R}^n$ 的单位球 $B_1$ 上的调和函数，则
\begin{align*}
H(r) = \fint_{\partial B_r} u^2\,\mathrm{d}S,\qquad D(r) = r^2\fint_{B_r} |\nabla u|^2\,\mathrm{d}y
\end{align*}
是 $r$ 的单调递增函数，其中 $B_r$ 是与 $B_1$ 同心，以 $r$ 为半径的球。
\end{theorem}
\begin{remark}
上述定理中的单调性在 $r=0$ 处也成立，但通常只对 $r>0$ 讨论。
\end{remark}
\begin{proof}
作伸缩变换，显然有
\begin{align*}
H(r) = \fint_{\partial B_1} u^2(rz)\,\mathrm{d}S(z).
\end{align*}
对 $r$ 求微商，由\hyperref[Basis of Analytics-theorem:含参量积分的可微性1]{含参量积分的可微性}和\hyperref[Basis of Analytics-theorem:Green第一恒等式]{Green第一恒等式}我们得到
\begin{align*}
H'(r) &= \fint_{\partial B_1} D u^2(rz)\cdot z\,\mathrm{d}S(z) = \fint_{\partial B_r} \frac{\partial u^2(y)}{\partial n}\,\mathrm{d}S(y) \\
&= \frac{1}{n\alpha(n)r^{n-1}}\int_{B_r} \Delta u^2\,\mathrm{d}y \xlongequal[\text{\rrefthe{Basis of Analytics-theorem:散度旋度Laplace算子的性质}{Basis of Analytics-theorem:散度旋度Laplace算子的性质-10}}]{u\text{是调和函数}} \frac{2}{n\alpha(n)r^{n-1}}\int_{B_r} |\nabla u|^2\,\mathrm{d}y \\
&= \frac{2}{nr}D(r)\geqslant 0,
\end{align*}
其中 $\bm{n}$ 为球面 $\partial B_r$ 的单位外法向量。

作伸缩变换，我们有
\begin{align}\label{eq::jiogjrJIOEFJi222}
D(r) = r^2\fint_{B_1} |\nabla u|^2(rz)\,\mathrm{d}z.
\end{align}
注意到
\begin{align*}
\frac{\partial}{\partial y_i}\left( u_iu_jy_j \right) =\left( u_iu_jy_j \right) _i=u_{ii}u_jy_j+u_iu_{ji}y_j+u_iu_j\delta _{ij}​\Longrightarrow u_iu_{ji}y_j=\left( u_iu_jy_j \right) _i-u_{ii}u_jy_j-u_iu_j\delta _{ij}.
\end{align*}
因为$u$调和,即$\sum_{i=1}^n{u_{ii}}=0$,所以$\sum_{i,j=1}^n{u_{ii}u_jy_j}=0$.故
\begin{align}\label{eq:jsfiojewiofjoiweijfii333-1}
\sum_{i,j=1}^n{u_iu_{ji}y_j}=\sum_{i,j=1}^n{\left[ \left( u_iu_jy_j \right) _i-u_iu_j\delta _{ij} \right]}.
\end{align}
并且
\begin{align}
\frac{2}{\alpha \left( n \right) r^{n-1}}\int_{B_r}{\left( \sum_{i,j=1}^n{u_iu_j\delta _{ij}} \right) \,\mathrm{d}y}=\frac{2r}{\alpha \left( n \right) r^n}\int_{B_r}{\left( \sum_{i=1}^n{u_{i}^{2}} \right) \,\mathrm{d}y}=\frac{2r}{\alpha \left( n \right) r^n}\int_{B_r}{\left| \nabla u \right|^2\,\mathrm{d}y}=\frac{2}{r}D\left( r \right) .\label{eq:jsfiojewiofjoiweijfii333-2}
\end{align}
因此对\eqref{eq::jiogjrJIOEFJi222}式关于$r$ 求微商，再由\hyperref[Basis of Analytics-theorem:含参量积分的可微性1]{含参量积分的可微性}可得
\begin{align*}
D'(r) &= 2r\fint_{B_1} |\nabla u|^2(rz)\,\mathrm{d}z + r^2\fint_{B_1} D(|\nabla u|^2(rz))\cdot z\,\mathrm{d}z 
\\
&=\frac{2}{r}D(r)+r^2\fint_{B_1}{D\left( \sum_{i=1}^n{u_{i}^{2}}\left( rz \right) \right) \cdot z\,\mathrm{d}z}
\\
&= \frac{2}{r}D(r) + 2r^2\fint_{B_1}\sum_{i,j=1}^n u_i(rz)u_{ij}(rz)z_j\,\mathrm{d}z \\
&= \frac{2}{r}D(r) + 2r\fint_{B_r}\sum_{i,j=1}^n u_i(y)u_{ji}(y)y_j\,\mathrm{d}y \\
&\xlongequal{\text{\eqref{eq:jsfiojewiofjoiweijfii333-1}式}} \frac{2}{r}D(r) + \frac{2}{\alpha(n)r^{n-1}}\int_{B_r}\sum_{i,j=1}^n [(u_i u_j y_j)_i - u_i u_j \delta_{ij}]\,\mathrm{d}y \\
&\xlongequal[\text{\eqref{eq:jsfiojewiofjoiweijfii333-2}式}]{\text{\hyperref[Basis of Analytics-theorem:Green恒等式-3]{Green恒等式}}} \frac{2}{\alpha(n)r^{n-1}}\int_{\partial B_r}\sum_{i,j=1}^n (u_i u_j y_j)\frac{y_i}{r}\,\mathrm{d}S(y) \\
&= 2nr\fint_{\partial B_r}\sum_{i,j=1}^n \left(u_i\frac{y_i}{r}\right)\left(u_j\frac{y_j}{r}\right)\mathrm{d}S(y) \\
&= 2nr\fint_{\partial B_r}\left|\frac{\partial u}{\partial n}\right|^2\,\mathrm{d}S \geqslant 0.
\end{align*}
因此 $D(r)$ 是 $r$ 的单调递增函数。

进一步，我们证明 $f(r)=\frac{D(r)}{H(r)}$ 是 $r$ 的单调递增函数。由上面的结果可知
\begin{align*}
D(r) = \frac{nr}{2}H'(r),
\end{align*}
以及
\begin{align*}
H'(r) = \int_{\partial B_r}\frac{\partial u^2}{\partial n}\,\mathrm{d}S,\qquad D'(r)=2nr\int_{\partial B_r}\left|\frac{\partial u}{\partial n}\right|^2\,\mathrm{d}S,
\end{align*}
于是
\begin{align*}
f'(r) &= H^{-2}\big(D'(r)H(r) - H'(r)D(r)\big) \\
&= H^{-2}\left(2nr\int_{\partial B_r}\left|\frac{\partial u}{\partial n}\right|^2\,\mathrm{d}S \cdot \int_{\partial B_r}u^2\,\mathrm{d}S - \frac{nr}{2}\left[\int_{\partial B_r}\frac{\partial u^2}{\partial n}\,\mathrm{d}S\right]^2\right) \\
&= 2nrH^{-2}\left(\int_{\partial B_r}|u_n|^2\,\mathrm{d}S \cdot \int_{\partial B_r}u^2\,\mathrm{d}S - \left[\int_{\partial B_r}u u_n\,\mathrm{d}S\right]^2\right) 
\geqslant 0.
\end{align*}
由 Cauchy 不等式，最后一个不等式成立。因此 $f(r)$ 是 $r$ 的单调递增函数.
\end{proof}

\begin{theorem}[调和函数比较定理]
设 \(\Omega\) 是 \(\mathbb{R}^n\) 中的有界开区域, \(u, v \in C^2(\Omega) \cap C(\overline{\Omega})\) 是调和函数. 若在边界 \(\partial\Omega\) 上成立 \(u(x) \leqslant v(x)\), 则在区域 \(\Omega\) 内成立 \(u(x) \leqslant v(x)\).
\end{theorem}

\begin{theorem}[Harnack第一定理]
如果 $\{u_k\}$ 中诸项都是有界区域 $\Omega$ 上的调和函数, 在 $\overline{\Omega}$ 上连续, 且该函数序列在 $\partial\Omega$ 上一致收敛, 那么该调和函数列在区域 $\Omega$ 也一致收敛, 且极限函数 $u$ 在区域 $\Omega$ 内也是调和函数.
\end{theorem}

\begin{theorem}[Harnack第二定理]
设 $\{u_k\}$ 是有界区域 $\Omega$ 上单调不减调和函数列, 若它在 $\Omega$ 内某一点 $P$ 收敛, 则它在 $\Omega$ 内处处收敛, 且极限函数是调和函数, 且收敛在任意闭子区域内是一致的.
\end{theorem}





\end{document}